\input{papers/template/style.tex}

\begin{document}
  \title{Causal Macroeconomic Modeling with Bayesian and Probabilistic
    Structural Models}

  \author{ADIA Authors}
  % TODO(*): Fill in the affiliation.
  \address{\textless{}Department\textgreater{},
    \textless{}Organization\textgreater{},
    \textless{}City, State/Country\textgreater{}}
  \email{@adia.ae}

  \keywords{causal inference, macroeconomics, Bayesian modeling,
    probabilistic modeling, structural causal models, point-in-time data}

  \date{\today}

  \begin{abstract}
    Macroeconomic variables enter causal models of asset returns in three
    roles: as common causes of firm characteristics and returns, as effect
    modifiers, and as priced risks. Using them raises problems that
    cross-sectional analysis avoids: one short non-randomized history,
    release lags and data revisions, mixed sampling frequencies, and regime
    changes.

    This paper formulates causal macroeconomic modeling as a probabilistic
    inference problem. A two-level structural causal model (SCM) is written
    as a hierarchical Bayesian generative model, with a macro layer indexed
    by time and a firm layer indexed by firm and time. Identification relies
    on externally identified shocks and on point-in-time data adapters.
    Posterior distributions over effects, regimes, and parameters make the
    weak identifying information of a short history explicit. The model
    extends the Theory Spec of a companion paper that converts finance papers
    into executable causal models.

    The paper presents the design and an evaluation protocol. No empirical
    evaluation has been carried out yet.
  \end{abstract}

  \maketitle

  \setcounter{tocdepth}{2}
  \tableofcontents

  % Status: OUTLINE stage. Each section lists the content it will contain.
  % Bullets are replaced by prose after the outline is approved.
  %
  % Provenance: the macro content (roles, time model, point-in-time data,
  % identification, estimator templates, literature-level graph) was moved
  % from `papers/Causal_Analysis_for_Finance` (see `notes.md`). The
  % probabilistic and Bayesian formulation (Sections~\ref{sec:model},
  % the last column of Table~\ref{tab:estimator_templates}, the last row of
  % Table~\ref{tab:macro_fields}, and the calibration and prior-sensitivity
  % checks) is new and does not come from `notes.md`.

  \section{Introduction}\label{sec:introduction}

  \begin{itemize}
  \tightlist
  \item
    Problem: macroeconomic variables are part of the causal structure of
    asset returns, but one short history makes them hard to model.

    \begin{itemize}
    \tightlist
    \item
      A cross-sectional graph can be static. Macro variables force a
      time-unrolled graph with lags and feedback.
    \item
      Omitting a macro common cause misattributes effects: a rate rise
      lowers valuations, which raises book-to-market, and also lowers
      realized returns directly. The value effect is then misattributed.
      This is a source of ``factor mirages''
      \cite{lopezdeprado2025primer}.
    \item
      There is one macro history, a few hundred monthly observations, and
      nothing is randomized. Point estimates and asymptotic standard errors
      hide how little the data identify.
    \end{itemize}
  \item
    Why a probabilistic formulation:

    \begin{itemize}
    \tightlist
    \item
      Short samples: parameter uncertainty is large and enters every
      downstream claim. A posterior carries it explicitly.
    \item
      Latent regimes: a posterior over regime paths replaces a hard
      classification of dates \cite{hamilton1989new}.
    \item
      Hierarchy: partial pooling of firm-level exposures to macro shocks
      \cite{gelman2007data}.
    \item
      Priors: shrinkage in vector autoregressions (VARs)
      \cite{doan1984forecasting,giannone2015prior}.
    \end{itemize}
  \item
    Scope relative to the companion paper \cite{adia2026finance}: the
    pipeline, the Theory Spec, and the factor mirage diagnostic are defined
    there and not repeated here. This paper defines only what macro
    variables add.
  \item
    Simplifying assumptions (bullet list, bold lead-in each):

    \begin{itemize}
    \tightlist
    \item
      \textbf{Macro layer exogenous to firms.} Edges run from macro to
      firm, not the reverse.
    \item
      \textbf{Public macro data only.} FRED-MD, FRED-QD, and ALFRED.
    \item
      \textbf{Linear-Gaussian components first.} Nonlinear and
      heavy-tailed variants are deferred.
    \end{itemize}
  \item
    Contributions, each tied to a section:

    \begin{itemize}
    \tightlist
    \item
      A two-level hierarchical probabilistic model with three macro roles
      (Section~\ref{sec:model}).
    \item
      Point-in-time data handling in the model's information set
      (Section~\ref{sec:point_in_time}).
    \item
      Identification with shocks and posterior uncertainty under weak
      identification (Section~\ref{sec:identification_with_shocks}).
    \item
      Estimator templates with probabilistic counterparts, and a simulator
      (Section~\ref{sec:estimators}).
    \item
      Theory Spec extensions and a macro-connected literature-level graph
      (Section~\ref{sec:literature_graph}).
    \item
      An evaluation protocol (Section~\ref{sec:evaluation}).
    \end{itemize}
  \item
    Roadmap: Section~\ref{sec:related_work} reviews related work and
    Section~\ref{sec:problem} formulates the problems.
    Section~\ref{sec:model} defines the model,
    Section~\ref{sec:point_in_time} the data layer, and
    Section~\ref{sec:identification_with_shocks} identification.
    Section~\ref{sec:estimators} covers estimators and the simulator, and
    Section~\ref{sec:literature_graph} the literature-level graph.
    Section~\ref{sec:evaluation} presents the evaluation protocol.
    Section~\ref{sec:limitations} discusses limitations and
    Section~\ref{sec:conclusion} concludes.
  \end{itemize}

  \section{Related Work}\label{sec:related_work}

  \begin{itemize}
  \tightlist
  \item
    \textbf{Macro variables in asset pricing.} Macro shocks as priced
    risks \cite{chen1986economic}; risk premia conditional on macro state
    \cite{ferson1991variation}. Contrast: each paper specifies one model.
    The role of the macro variable (confounder, modifier, priced risk) is
    left implicit. This paper makes the role an explicit edge attribute.
  \item
    \textbf{Identified macro shocks.} Monetary policy surprises
    \cite{gertler2015monetary}, narrative shocks
    \cite{romer2004new}, oil supply shocks \cite{kilian2009oil}; dynamic
    effects by local projections \cite{jorda2005local}. Contrast: these
    supply the shocks and the estimator. This paper makes the shock a
    node type that the validator checks.
  \item
    \textbf{Time-series causal discovery.} PCMCI \cite{runge2019pcmci}
    and VAR-LiNGAM \cite{hyvarinen2010var}. Contrast: on short samples
    the output is noisy, so it serves as a consistency check only.
  \item
    \textbf{Bayesian macroeconometrics.} VARs \cite{sims1980macroeconomics},
    shrinkage priors \cite{doan1984forecasting,giannone2015prior},
    time-varying parameter VARs with stochastic volatility
    \cite{primiceri2005time}, Bayesian dynamic stochastic general
    equilibrium (DSGE) models \cite{smets2007shocks}, regime switching
    \cite{hamilton1989new}. Contrast: these target macro forecasting and
    policy analysis. Here they form the macro layer of a model with a firm
    layer below it.
  \item
    \textbf{Real-time macro data.} The real-time data set of
    \cite{croushore2001real} and the FRED-MD database
    \cite{mccracken2016fredmd}. Contrast: these are data resources. Here
    vintage and release lag are node attributes of the Theory Spec.
  \end{itemize}

  \section{Problem Formulation}\label{sec:problem}

  \subsection{Notation and Setup}\label{sec:notation_and_setup}

  \begin{itemize}
  \tightlist
  \item
    A Theory Spec \(\mathcal{T}\) as defined in the companion paper
    \cite{adia2026finance}, with the macro extensions of
    Section~\ref{sec:literature_graph}.
  \item
    Macro layer: a vector of macro series \(\mathbf{m}_t\) for
    \(t = 1, \dots, T\) and a latent regime \(s_t\).
  \item
    Firm layer: characteristics \(x_{i,t}\) and returns \(r_{i,t+1}\) for
    firm \(i\).
  \item
    Parameters \(\theta = (\theta_m, \theta_f)\) for the macro and firm
    layers.
  \item
    Information set \(\mathcal{D}_d\): the data vintage that was knowable
    on date \(d\).
  \item
    Target: the posterior \(p(\theta, s_{1:T} | \mathcal{D}_d,
    \mathcal{T})\) and the induced posterior of the estimand \(\tau\).
  \end{itemize}

  \subsection{Simplifying Assumptions}\label{sec:simplifying_assumptions}

  \begin{itemize}
  \tightlist
  \item
    Bullet list, bold lead-in each: the macro layer is exogenous to
    firms; time-indexing removes feedback cycles; macro names come from a
    small canonical list (real policy rate, term spread, credit spread,
    industrial production growth); components are linear-Gaussian at
    first.
  \end{itemize}

  \subsection{Problem 1: Role Assignment}\label{sec:problem_1_role_assignment}

  \begin{itemize}
  \tightlist
  \item
    For every edge from a macro node, decide whether it means common
    cause, effect modifier, or direct cause (priced risk), and record the
    role in \(\mathcal{T}\).
  \item
    Quality measured by agreement with expert labels on a set of papers.
  \end{itemize}

  \subsection{Problem 2: Point-in-Time Information
  Set}\label{sec:problem_2_point_in_time}

  \begin{itemize}
  \tightlist
  \item
    Given \(\mathcal{T}\) and a date \(d\), build \(\mathcal{D}_d\) such
    that every observation was published on or before \(d\), under the
    vintage policy that the paper's theory requires.
  \end{itemize}

  \subsection{Problem 3: Identification and Posterior
  Inference}\label{sec:problem_3_identification_and_inference}

  \begin{itemize}
  \tightlist
  \item
    Decide whether \(\tau\) is identified from an exogenous shock or
    from the endogenous policy rate.
  \item
    Compute the posterior of \(\tau\) given \(\mathcal{D}_d\), integrating
    over regime and parameter uncertainty.
  \item
    Report the prior next to the posterior, so that the identifying
    content of the data is visible.
  \end{itemize}

  \section{Probabilistic Two-Level Model}\label{sec:model}

  \subsection{Three Roles of Macro
  Nodes}\label{sec:three_roles_of_macro_nodes}

  \begin{itemize}
  \tightlist
  \item
    \textbf{Common cause.} Rates, credit spreads, and the business cycle
    drive both characteristics and returns. Omission yields misattributed
    effects. Probabilistic reading: the macro state enters both the
    characteristic and the return equations, so omitting it biases the
    posterior of the characteristic effect.
  \item
    \textbf{Effect modifier.} Regime-dependent claims, such as momentum
    crashing after market rebounds or value paying off in recoveries
    (\texttt{modulated\_by} edge attribute or interaction node).
    Probabilistic reading: the coefficient depends on the regime \(s_t\).
  \item
    \textbf{Direct cause or priced risk.} Macro shock with firm-level
    exposure (beta) \cite{chen1986economic}. Probabilistic reading: the
    exposures \(\beta_i\) are draws from a common distribution.
  \end{itemize}

  \subsection{Two-Level Hierarchical Model}\label{sec:two_level_model}

  \begin{itemize}
  \tightlist
  \item
    Two-level SCM \cite{pearl2009causality}: the macro layer is indexed by
    time, the firm layer by firm and time. Edges run from macro to firm,
    not the reverse.
  \item
    Generative description:

    \begin{itemize}
    \tightlist
    \item
      Regime: \(s_t | s_{t-1} \sim \text{Markov}(\Pi)\).
    \item
      Macro layer: \(\mathbf{m}_t | \mathbf{m}_{t-1}, \dots,
      \mathbf{m}_{t-p}, s_t \sim p(\cdot | \theta_m)\), e.g., a VAR with
      regime-dependent coefficients.
    \item
      Firm layer: \(r_{i,t+1} | x_{i,t}, \mathbf{m}_t, s_t \sim
      p(\cdot | \beta_i, \theta_f)\), with \(\beta_i\) drawn from a common
      prior (partial pooling).
    \end{itemize}
  \item
    Figure~\ref{fig:two_level_scm} shows the two-level SCM.
  \end{itemize}

  \begin{figure}[H]
    \centering
    % TODO(*): Replace the placeholder with the two-level SCM diagram.
    \fbox{\parbox{0.9\linewidth}{\centering\vspace{3em}
      Figure placeholder\vspace{3em}}}
    \caption{Two-level structural causal model with a macro layer and a firm
      layer.}
    \label{fig:two_level_scm}
  \end{figure}

  \subsection{Time Model}\label{sec:time_model}

  \begin{itemize}
  \tightlist
  \item
    Lags on every edge, contemporaneous flag, feedback made acyclic by
    unrolling.
  \item
    Mixed frequency: native frequency per node, aggregation rule per edge
    (end-of-period, average, MIDAS weights).
  \item
    Option: treat a low-frequency series as a partial observation of a
    high-frequency latent process.
  \item
    Stationarity transforms recorded per series (FRED-MD codes
    \cite{mccracken2016fredmd}). Levels of trending series in a regression
    are the classic spurious-regression trap.
  \end{itemize}

  \subsection{Priors and Regimes}\label{sec:priors_and_regimes}

  \begin{itemize}
  \tightlist
  \item
    Shrinkage priors on VAR coefficients, with the prior tightness
    estimated or set by hierarchical choice
    \cite{doan1984forecasting,giannone2015prior}.
  \item
    Markov-switching regimes with a transition matrix \(\Pi\); the
    posterior over \(s_{1:T}\) is an output.
  \item
    Partial pooling of firm exposures \(\beta_i\) toward a shared mean.
  \item
    Prior sensitivity: report the posterior of \(\tau\) under at least two
    prior choices.
  \item
    Inference by Markov chain Monte Carlo (MCMC) or variational methods in
    a probabilistic programming library (e.g., PyMC or Stan).
  \end{itemize}

  % TODO(*): Decide the prior families and the inference engine.

  \section{Point-in-Time Data}\label{sec:point_in_time}

  \begin{itemize}
  \tightlist
  \item
    Revisions and release lags make revised data a source of look-ahead
    bias. GDP for the first quarter is first published weeks later and then
    revised for years. Generated code uses today's revised numbers by
    default.
  \item
    Fields: \texttt{release\_lag}, the time from the period's end to first
    publication, and \texttt{vintage\_policy} (first release, real-time
    vintage, final revised), chosen by what the paper's theory requires.
  \item
    Adapter signature \texttt{(series,\ as\_of\_date)}, returning only what
    was knowable on that date. Sources: ALFRED, the Philadelphia Fed
    Real-Time Data Set \cite{croushore2001real}, FRED-MD and FRED-QD
    \cite{mccracken2016fredmd}.
  \item
    Probabilistic reading: the information set \(\mathcal{D}_d\) is the
    conditioning set of the posterior.
  \item
    Open question: model a first release as a noisy measurement of the
    final value, with the revision variance estimated from the vintages.
  \end{itemize}

  \section{Identification with Shocks}\label{sec:identification_with_shocks}

  \begin{itemize}
  \tightlist
  \item
    One short, non-randomized macro history limits identification.
  \item
    \texttt{instrument} node type for identified shocks: monetary
    surprises \cite{gertler2015monetary}, narrative shocks
    \cite{romer2004new}, oil supply shocks \cite{kilian2009oil}.
  \item
    Validator check: does the estimate use an identified shock or the
    endogenous policy rate?
  \item
    Weak identification: when the shock is weak, the posterior of
    \(\tau\) is expected to stay close to the prior. Report both.
  \item
    Time-series discovery as a consistency check only: PCMCI
    \cite{runge2019pcmci}, VAR-LiNGAM \cite{hyvarinen2010var}, Granger
    as a weak baseline.
  \item
    Structural break and regime tests (Bai-Perron \cite{bai1998estimating},
    Hamilton \cite{hamilton1989new}). A macro edge that holds before 2008
    and vanishes after is a finding, not a bug.
  \end{itemize}

  \section{Estimators and Simulator}\label{sec:estimators}

  \begin{itemize}
  \tightlist
  \item
    Table~\ref{tab:estimator_templates} maps each paper claim to an
    estimator template and to its probabilistic counterpart.
  \item
    Simulator: macro layer (VAR or regime switching) first, firm layer
    conditional on it, to test estimator recovery under macro confounding.
    The simulator is the generative model of Section~\ref{sec:model} run
    forward with known parameters.
  \item
    Posterior predictive checks compare simulated and observed macro
    series.
  \end{itemize}

  \begin{table}[H]
    \centering
    \begin{tabularx}{\linewidth}{lXX}
      \toprule
      Paper's claim & Estimator template & Probabilistic counterpart \\
      \midrule
      Dynamic effect of a macro shock & Local projections
      \cite{jorda2005local}, SVAR (\texttt{statsmodels}) & Bayesian VAR with a
      shrinkage prior \cite{giannone2015prior}; posterior impulse responses
      \\
      Macro-conditional factor premium & Conditional Fama-MacBeth with macro
      interactions \cite{ferson1991variation} & Hierarchical regression with
      time-varying coefficients \cite{primiceri2005time} \\
      Regime-dependent effect & Markov-switching regression
      \cite{hamilton1989new} & Bayesian Markov-switching regression;
      posterior regime probabilities \\
      Macro risk priced in the cross-section & Two-pass beta estimation on
      macro shock series & Joint hierarchical model of exposures and risk
      premium \\
      Mixed-frequency prediction & MIDAS regression & Bayesian MIDAS or
      mixed-frequency state-space model \\
      \bottomrule
    \end{tabularx}
    \caption{Estimator templates for macro claims and their probabilistic
      counterparts.}
    \label{tab:estimator_templates}
  \end{table}

  \section{Literature-Level Graph}\label{sec:literature_graph}

  \begin{itemize}
  \tightlist
  \item
    Macro is the glue between graphs. Canonical macro names (real policy
    rate, term spread, credit spread, industrial production growth)
    connect graphs from macro papers and from cross-sectional finance
    papers into pathways that no single paper tested, e.g., policy to
    credit spread to distress returns.
  \item
    The system can generate the combined model and test it.
  \item
    Table~\ref{tab:macro_fields} lists the Theory Spec fields added for
    macro nodes, edges, and validation.
  \end{itemize}

  \begin{table}[H]
    \centering
    \begin{tabularx}{\linewidth}{lX}
      \toprule
      Component & Added fields or checks \\
      \midrule
      Node & \texttt{level} (macro, firm, or industry), \texttt{frequency},
      \texttt{transform}, \texttt{release\_lag}, \texttt{vintage\_policy},
      \texttt{source\_series\_id} \\
      Edge & \texttt{lag}, \texttt{contemporaneous} flag,
      \texttt{modulated\_by}, \texttt{aggregation\_rule},
      \texttt{regime\_scope} \\
      New node type & \texttt{instrument} or shock, with the source study for
      its identification \\
      Validation & Time-series discovery check, endogeneity check of the macro
      regressor, structural-break check \\
      Probabilistic model & \texttt{prior}, \texttt{pooling} (none, partial,
      or complete), \texttt{regime\_variable} \\
      \bottomrule
    \end{tabularx}
    \caption{Theory Spec extensions for macroeconomic variables.}
    \label{tab:macro_fields}
  \end{table}
  % TODO(*): Decide the names of the probabilistic model fields.

  \section{Evaluation Protocol}\label{sec:evaluation}

  \begin{itemize}
  \tightlist
  \item
    Status: protocol only. No results are reported.
  \item
    First milestone: the macro confounder role, a handful of FRED-MD series
    with point-in-time handling, applied to the cross-sectional anomaly
    papers of the companion paper \cite{adia2026finance} and the signals of
    \cite{chen2022open}.
  \item
    \textbf{Role assignment.} Agreement of the extracted macro roles with
    expert labels (Problem~1).
  \item
    \textbf{Factor survival.} Which published factors survive macro
    adjustment.
  \item
    \textbf{Recovery.} In the simulator, does the estimator recover known
    effects when macro confounding is present?
  \item
    \textbf{Calibration.} In the simulator, do posterior credible intervals
    cover the true effect at the nominal rate?
  \item
    \textbf{Vintage ablation.} Run the same estimates under first-release,
    real-time, and final-revised vintages. Measure the look-ahead effect.
  \item
    Baselines: the same claim without the macro node; the frequentist
    estimator template of Table~\ref{tab:estimator_templates}.
  \item
    Table~\ref{tab:evaluation} lists metrics, data, and baselines.
  \end{itemize}

  \begin{table}[H]
    \centering
    \begin{tabularx}{\linewidth}{lXXX}
      \toprule
      Component & Metrics & Data & Baseline or variant \\
      \midrule
      Role assignment & Agreement with expert labels (accuracy per edge) &
      Anomaly papers with expert-labeled macro edges & Not applicable \\
      Factor survival & Share of published factors with a significant effect
      after macro adjustment & Anomaly papers, FRED-MD series & The same claim
      without the macro node \\
      Recovery & Bias and error of the estimated effect & Simulator output &
      Frequentist estimator template \\
      Calibration & Coverage of credible intervals & Simulator output & Not
      applicable \\
      Vintage ablation & Change in the estimate across vintage policies &
      ALFRED vintages & Final-revised data \\
      \bottomrule
    \end{tabularx}
    \caption{Evaluation protocol: components, metrics, data, and baselines.}
    \label{tab:evaluation}
  \end{table}
  % TODO(*): Replace with a Technical Evaluation section once an
  % implementation exists.

  \section{Discussion and Limitations}\label{sec:limitations}

  \begin{itemize}
  \tightlist
  \item
    \textbf{Look-ahead bias.} Release lags and data revisions are the main
    failure points of generated code.
  \item
    \textbf{Small samples.} Macro history is short, so discovery
    algorithms are noisy and serve only as consistency checks. Posteriors
    are expected to be wide.
  \item
    \textbf{Prior sensitivity.} With little data, conclusions can depend
    on the prior. Report results under more than one prior.
  \item
    \textbf{Regime instability.} Estimated regimes and breaks can change
    with the sample end date.
  \item
    \textbf{Computation.} Hierarchical models over many firms need
    scalable inference.
  \item
    \textbf{Endogeneity.} Without an identified shock, a macro regressor
    such as the policy rate is endogenous to the economy.
  \end{itemize}

  \section{Conclusion and Future Work}\label{sec:conclusion}

  \begin{itemize}
  \tightlist
  \item
    Summary of contributions in 2 to 3 sentences.
  \item
    State plainly: design and protocol only, no implementation or
    empirical evaluation yet.
  \item
    Numbered next steps:

    \begin{enumerate}
    \def\labelenumi{\arabic{enumi}.}
    \tightlist
    \item
      Extend the Theory Spec schema with the macro fields.
    \item
      Build the point-in-time adapter for ALFRED and FRED-MD.
    \item
      Implement the estimator templates and their probabilistic
      counterparts.
    \item
      Build the simulator and run the recovery and calibration checks.
    \item
      Add the macro confounder role to the anomaly-paper evaluation.
    \item
      Connect macro papers and finance papers in the literature-level
      graph.
    \end{enumerate}
  \end{itemize}

  \bibliographystyle{amsplain}
  \bibliography{references}
\end{document}
